% Сборка из корня проекта: tectonic docs/documentation.tex
% Альтернатива: xelatex -output-directory=docs docs/documentation.tex (дважды).
\documentclass[11pt,a4paper]{article}
\usepackage{fontspec}
\setmainfont{Liberation Serif}
\setsansfont{Liberation Sans}
\setmonofont{Liberation Mono}[Scale=0.85]
\usepackage[russian]{babel}
\usepackage[margin=20mm]{geometry}
\usepackage{amsmath,booktabs,tabularx,array,enumitem,tikz}
\usetikzlibrary{arrows.meta,positioning}
\usepackage[unicode,hidelinks]{hyperref}
\setlength{\parindent}{0pt}
\setlength{\parskip}{5pt}
\setlist{nosep,leftmargin=*}
\newcommand{\code}[1]{\texttt{#1}}
\newcolumntype{Y}{>{\raggedright\arraybackslash}X}
\begin{document}
\sloppy
\begin{center}
{\Large\bfseries Система управления инвестиционным портфелем}\\[5pt]
{\large Документация программной системы}\\[3pt]
Учебное имитационное моделирование · C++17 / Qt 6
\end{center}

\section{Назначение и постановка задачи}
Программа моделирует управление инвестиционным фондом: пользователь распределяет капитал между активами, наблюдает изменение рынка, учитывает доход, налог и операции с паями, затем изменяет состав портфеля. Цель эксперимента --- исследовать влияние структуры вложений, валютных курсов и потоков средств на капитал, прибыль и ликвидность фонда.

Модель работает локально, без подключения к бирже. Поддерживаются рублёвые и валютные вклады, акции, облигации и драгоценные металлы. Шаг моделирования фиксирован: один месяц; продолжительность партии --- от 12 до 30 месяцев. Перед началом задаются капитал, налоговая ставка, seed генератора, валютные курсы, диапазоны их изменений, условия инструментов и режим движения паёв.

Результаты: месячные отчёты, чистая прибыль, налоги, итоговый капитал, накопленная доходность, максимальная просадка, выпуск и погашение паёв. Визуализация включает графики, структуру активов и анимированное раскрытие графика при обновлении данных.

\subsection*{Метод моделирования и основные формулы}
Используется дискретная имитационная модель с псевдослучайными изменениями рынка. Генератор --- \code{std::mt19937\_64}; изменения цен и курсов равномерно распределены в заданных границах:
\[
 P_t=P_{t-1}(1+\varepsilon_t),\qquad \varepsilon_t\sim U(a,b).
\]
Чистая стоимость фонда (NAV) равна свободным рублям плюс стоимости позиций в рублях за вычетом обязательств. Для месяца с начальным капиталом $A_0$ и оценкой $A_1$ после изменения рынка, но до налога и потоков:
\[
 \Pi=A_1-A_0,\quad T=\operatorname{round}(\max(\Pi,0)\tau),\quad
 \Pi_{net}=\Pi-T,\quad r=\Pi_{net}/A_0.
\]
Поступления и выплаты пайщикам не являются инвестиционной прибылью. Доходность за несколько месяцев рассчитывается как $\prod_t(1+r_t)-1$. При неположительном начальном NAV сравнимая доходность не определяется.

В автоматическом режиме спрос зависит от результата месяца:
\[
 d=\operatorname{clamp}(kr+U(-n,n),-d_{out},d_{in}).
\]
Положительное $d$ задаёт поступление средств, отрицательное --- долю возвращаемых паёв. Запрос рассчитывается один раз и подтверждается пользователем.

\subsection*{Допущения и ограничения}
Проценты вклада простые, ставка фиксируется при открытии; досрочное закрытие запрещено. Облигация оценивается по номиналу с накопленным купоном. Акции не выплачивают дивиденды. Комиссии, спред и кредитное плечо отсутствуют. Налог начисляется и на нереализованную переоценку; убытки не переносятся. Это учебные правила, а не воспроизведение банковского или налогового учёта.

\newpage
\section{Классы и взаимодействие объектов}
\begin{center}
\begin{tikzpicture}[font=\small,box/.style={draw,rounded corners,align=center,minimum width=38mm,minimum height=12mm},dep/.style={-{Stealth},dashed},own/.style={Diamond-},inherit/.style={-{Triangle[open]}}]
\node[box] (mw) at (0,0) {MainWindow};
\node[box] (qm) at (0,1.8) {QMainWindow};
\node[box] (dw) at (5,0) {DashboardWidget};
\node[box] (qw) at (5,1.8) {QWidget};
\node[box] (sim) at (0,-2) {Simulation};
\node[box] (gs) at (5,-2) {GameStorage};
\node[box] (state) at (0,-4) {GameState};
\node[box] (port) at (-4.5,-2) {Portfolio};
\draw[inherit] (mw)--(qm);
\draw[inherit] (dw)--(qw);
\draw[own] (mw)--node[left]{1}(sim);
\draw[own] (mw)--node[above]{1}(dw);
\draw[dep] (mw)--(gs);
\draw[own] (sim)--node[left]{1}(state);
\draw[dep] (sim)--(port);
\draw[dep] (gs)--(state);
\draw[dep] (dw.east) -- ++(0.6,0) |- (state.east);
\end{tikzpicture}
\end{center}
Пустой треугольник обозначает наследование, ромб --- владение, пунктирная стрелка --- использование. Диаграмма показывает основные классы; служебные Qt-классы и внутренние виджеты графиков опущены. Иерархия наследования используется в интерфейсе; виды активов представлены данными и \code{InstrumentKind}, а не отдельной иерархией подклассов.

В окне находится один объект \code{Simulation}, владеющий состоянием партии и генератором. \code{Portfolio} предоставляет статические операции над переданным портфелем. \code{DashboardWidget} хранит данные отображения, а команды передаёт окну сигналами. Расчётные классы не зависят от Qt. Все действия выполняются в GUI-потоке.

\subsection*{Спецификации основных классов}
\begin{description}[style=nextline]
\item[Simulation] \code{start(config)} проверяет настройки и начинает игру; \code{trade(request)} выполняет сделку; \code{previewTrade(request)} возвращает предварительный расчёт без изменения игры. \code{confirmInitialAllocation()}, \code{nextMonth()}, \code{applyFlows(request)}, \code{finishMonth()} управляют фазами. \code{state()} даёт доступ только для чтения, \code{statistics()} --- итоги. \code{exportState()} и \code{restoreState(candidate)} обеспечивают сохранение и восстановление. Ошибка команды не меняет состояние и генератор.
\item[Portfolio] \code{executeTrade(...)} изменяет деньги и позиции; \code{valuePosition(...)} и \code{assetsValue(...)} оценивают активы; \code{accrueAndRedeem(...)} начисляет фиксированный доход и погашает инструменты. \code{validate(...)} проверяет портфель. Некорректная сделка отклоняется без частичного выполнения.
\item[GameStorage] \code{save(path,state)} атомарно записывает JSON; \code{load(path)} возвращает проверенное состояние-кандидат. Ошибка чтения не заменяет активную игру. Формат хранит фазу, историю, незакрытый месяц и состояние генератора.
\item[MainWindow и DashboardWidget] Окно собирает ввод, вызывает модель и обновляет представление. Метод обзора \code{setState(state,stats)} обновляет графики и карточки. Сигналы \code{monthRequested(int)} и \code{positionRequested(qint64)} открывают отчёт или действие по устойчивому идентификатору.
\end{description}
Общие результаты команд представлены \code{Result<T>}: успешное значение либо код и сообщение ошибки. Деньги хранятся в целых копейках/центах, паи --- с точностью $10^{-6}$, металл --- $0{,}001$ грамма.

\newpage
\section{Инструментальные средства и файлы}
Использованы C++17, стандартная библиотека C++, Qt 6 Core/Widgets/Test, Qt Designer (формат \code{.ui}), CMake и GCC. Графики рисуются через QPainter; анимация использует QVariantAnimation. Проект можно открыть в Qt Creator через \code{CMakeLists.txt}; конкретная IDE не требуется для сборки. Проверенное окружение: Linux, GCC 16.2.1, Qt 6.12.0. Из-за применения \code{\_\_int128} для Windows нужен комплект MinGW, поддержка MSVC не заявлена.

\begin{tabularx}{\linewidth}{@{}p{64mm}Y@{}}
\toprule
Файл / группа файлов & Назначение \\\midrule
\code{CMakeLists.txt} & Цели сборки приложения, библиотек и тестов.\\
\code{src/main.cpp} & QApplication, локаль, запуск окна.\\
\code{src/GameTypes.h} & Настройки, инструменты, позиции, фазы, отчёты, состояние и проверяемая арифметика.\\
\code{src/Simulation.h/.cpp} & Месячный цикл, рынок, налог, паи, ликвидность и статистика.\\
\code{src/Portfolio.h/.cpp} & Сделки, оценка, проценты, купоны и погашения.\\
\code{src/GameStorage.h/.cpp} & Сериализация и проверка JSON.\\
\code{src/MainWindow.h/.cpp/.ui} & Формы, меню, вкладки, таблицы и обработка команд.\\
\code{src/DashboardWidget.h/.cpp} & Обзор, графики, структура активов, карточки и шкала времени.\\
\code{tests/SimulationTests.cpp} & Проверки модели, сохранений и пользовательских сценариев.\\
\code{README.md}, \code{AGENTS.md} & Запуск и правила сопровождения проекта.\\
\code{docs/project-plan.md} & Карта проекта и взаимодействие модулей.\\
\code{docs/simulation-rules.md} & Подробные формулы и принятые допущения.\\
\code{docs/documentation.tex/.pdf} & Исходник и собранная версия этой документации.\\
\bottomrule
\end{tabularx}

\subsection*{Сборка и запуск}
Необходимы CMake, GCC либо Clang, Qt 6 Widgets и Test. Из корневой папки проекта:
\begin{verbatim}
cmake -S . -B build -DCMAKE_BUILD_TYPE=Debug
cmake --build build --parallel 2
./build/portfolio_simulator
ctest --test-dir build --output-on-failure
\end{verbatim}
Тесты интерфейса выполняются в режиме \code{offscreen}. Для сборки без тестов передаётся \code{-DBUILD\_TESTING=OFF}. PDF собирается командой \code{tectonic docs/documentation.tex}; нужны шрифты Liberation. Возможна сборка XeLaTeX (два прохода, каталог вывода \code{docs}).

\subsection*{Хранение состояния и репозиторий}
Сохранение --- JSON с \code{schemaVersion: 1}. Целые денежные величины, количества и идентификаторы записываются строками, чтобы избежать потери точности. Предельный размер файла --- 10 МиБ. Загрузка проверяет типы, диапазоны, связи и согласованность отчётов. Повторяемость случайной последовательности гарантируется в том же окружении стандартной библиотеки.

Репозиторий проекта:\\
\url{https://github.com/BlackDiamond2108/investment_portfolio_simulation}

\newpage
\section{Интерфейс и порядок работы}
Главное окно содержит постоянную панель вкладок и действия сохранения. На узком экране настройки перестраиваются в одну колонку, длинные страницы прокручиваются.

\begin{description}[style=nextline]
\item[Обзор] Баланс фонда, свободные средства, обязательства, прибыль и просадка; график капитала или месячной прибыли, кольцевая диаграмма структуры активов, карточки вложений и шкала месяцев. Период графика переключается между всей историей, 6 и 12 месяцами. Наведение показывает точные значения, нажатие на точку открывает отчёт.
\item[Рынок] Карточки доступных предложений и подробная таблица условий; переход к покупке выбранного инструмента.
\item[Операции] Покупка и продажа активов, открытие вкладов, предварительный расчёт сделки, ввод и подтверждение поступлений и возврата паёв.
\item[История и итоги] Месячные суммы капитала, дохода, налога и потоков; детализация дохода по позициям. Предварительный месяц отмечен отдельно от закрытых.
\item[Новая игра] Капитал, число месяцев, налог, seed, валюты, диапазоны изменения рынка и параметры спроса. Редактор каталога позволяет менять и добавлять инструменты. Изменение настроек не изменяет уже начатую партию.
\end{description}

\subsection*{Месячный цикл}
\begin{enumerate}
\item Создать игру и собрать начальный портфель, оставив при необходимости часть капитала свободной. Подтвердить распределение.
\item Нажать «Рассчитать следующий месяц». Модель изменит рынок, начислит доход, выполнит погашения и рассчитает налог.
\item Изучить результаты и подтвердить потоки: сумму поступлений и число возвращаемых паёв либо предложение автоматической модели. Допускаются нулевые потоки.
\item Изменить портфель. При наличии обязательств продать ликвидные активы: выручка сначала покрывает налог, затем долг пайщикам.
\item Нажать «Завершить месяц». После последнего месяца получить итоги; в последнем месяце новые покупки запрещены.
\end{enumerate}
Фазы модели: начальное распределение, готовность к месяцу, ожидание потоков, реструктуризация, завершение. Повторное нажатие не приводит к повторному начислению. При невозможности рассчитаться по долгам игра завершается досрочно; обязательства остаются отражены в состоянии.

\subsection*{Вклады, графики и сохранение}
На карточке вклада кнопка «Условия» показывает ставку и месяц автоматического закрытия. По наступлении срока сумма с процентами переводится в свободные рубли по текущему курсу. Вклад, срок которого выходит за пределы игры, остаётся в итоговой оценке с начисленными процентами.

Движущийся элемент интерфейса --- анимированное раскрытие графика длительностью 400 мс. Анимация меняет только отображение: экономическое время продвигается по команде пользователя. Рост графика капитала может быть вызван поступлениями; инвестиционный результат следует смотреть на графике прибыли. Доли кольцевой диаграммы считаются по активам до вычета обязательств.

«Открыть пример» создаёт настоящую партию с seed=42, смешанными вложениями и шестью рассчитанными месяцами. Её можно продолжить. Сохранение и открытие доступны в меню «Партия» и по Ctrl+S / Ctrl+O. При замене несохранённой игры предлагается сохранить изменения.

\newpage
\section{Контрольные эксперименты и проверка}
Эксперименты выполнены на расчётном движке через сценарии \code{SimulationTests}. Это небольшие контролируемые прогоны для проверки зависимостей, а не статистическая оценка выгодности стратегий. Случайность в контрольных случаях отключена нулевыми диапазонами либо заменена фиксированным изменением. Ниже суммы приведены в рублях, если не указана другая валюта.

\subsection*{1. Процентный доход и налог}
Сценарий \code{fixedDeposit}: капитал 12\,000, вся сумма размещена в рублёвый вклад на 12 месяцев под 12\% годовых, налог 17\%. После первого шага доход составил 120,00, налог --- 20,40, чистая прибыль --- 99,60. Результат подтверждает ежемесячное начисление простого процента и удержание налога только с положительной прибыли. Это проверка первого месяца, а не успешного завершения всей партии: отсутствие свободных денег может потребовать расчёта по обязательствам.

\subsection*{2. Изменение валютного курса}
В сценарии \code{market} при нулевых диапазонах цены и курс остались неизменными. Во втором варианте открыт вклад 100 USD с нулевой ставкой и задан рост USD на 10\% за месяц: курс изменился с 90 до 99 RUB/USD, а стоимость вклада --- с 9\,000 до 9\,900 RUB. Даже без процентов валютная позиция даёт рублёвую переоценку.

Дополнительный прогон \code{bondAndCurrencyIncome}: те же 100 USD, ставка 12\%, срок два месяца, налог 0\%, ежемесячный рост курса 10\%. Доход первого месяца --- 999,00, второго --- 1\,108,80; суммарно --- 2\,107,80. После второго месяца вклад погашен. Процентный доход и валютная переоценка учитываются совместно, без повторного включения погашаемой основной суммы в прибыль.

\subsection*{3. Поступления пайщиков и прибыль}
Сценарий \code{flows} повторяет рублёвый вклад первого опыта. После расчёта месяца вводится поступление в размере стоимости двух паёв. Цена пая составляет 1\,008,30; принято 2\,016,60. Капитал после налога увеличивается с 12\,099,60 до 14\,116,20, но прибыль от вложений остаётся 120,00 до налога. Эксперимент показывает, почему рост капитала нельзя отождествлять с доходностью.

\subsection*{4. Купон при покупке облигации между выплатами}
Сценарий \code{bondAndCurrencyIncome}: налог 0\%, ставка купона 12\%, погашение в месяце 8. В конце месяца 2 приобретаются пять облигаций номиналом 1\,000 за 5\,100 с учётом накопленного купона. В месяцах 3--8 прибыль составляет по 50,00; итог --- 300,00, после погашения позиций не остаётся. Уплаченный при покупке купон не становится дополнительной прибылью покупателя.

\subsection*{Проверки и воспроизведение}
Полный набор CTest выполнен успешно при подготовке документации. Он включает расчёты, ошибочные команды, переполнения, сохранение и продолжение генератора, прохождение 12 и 30 месяцев через окно, диаграммы и навигацию. Для повторения описанных опытов:
\begin{verbatim}
QT_QPA_PLATFORM=offscreen ./build/simulation_tests \
  fixedDeposit market flows bondAndCurrencyIncome
\end{verbatim}
\end{document}
